%% sections/limitations.tex
%% Honest Limitations & Ethics (ML-venue norm).

\section{Limitations and Ethics}
\label{sec:limits}

\paragraph{Scope of the generator.} The IR targets composable sequences of the four
operations in \Cref{tab:ops}, not arbitrary control flow, and
\texttt{AGGREGATE}/\texttt{CONDITIONAL\_SELECT} are restricted to integer operands. This
is a deliberate trade rather than an accident: integer-only results format
byte-identically across all five backends, which is what lets a single oracle check
equivalence across Python, JavaScript, Go, Java, and C++. Generated artifacts are
single-file programs, not repository-scale projects. Consequently the dataset is a
\emph{controlled stimulus}, and we claim findings about behavior on that stimulus ---
robustness to incidental complexity, equivalence-preserving transformation --- not that
absolute results transfer to real-world codebases. Scaling the instance count does not
widen this scope; broadening the operation set, the operand types, and the program
granularity is future work.

\paragraph{The readability leg of construct validity is weak.} The label order is
validated, not asserted (\Cref{sec:validity}), but the legs of that validation are
uneven. The complexity metrics move strongly with the knob, yet that is partly
\emph{by construction}: the same transforms that define a higher knob level (e.g.\ adding
a branch) mechanically inflate cyclomatic complexity, so those metrics are convergent but
not independent witnesses. The genuinely more independent leg is readability, and it is
weak: the surface-feature aggregate correlates with the knob at only $\approx0.39$
(\Cref{sec:validity}), and it is a faithful re-implementation of the published
Buse--Weimer \emph{features}~\citep{buse2010readability}, \emph{not} the original
human-calibrated trained model (whose coefficients we could not obtain offline), so we
report feature correlations rather than a calibrated readability score. A fresh,
controlled human ranking on these exact instances --- the one experiment that would fully
dissolve the residual circularity between the knob and the complexity metrics --- is
\emph{not} run here and is the principal open item; we mark it as future work and do not
overclaim the by-construction order as validated human quality.

\paragraph{Contamination is mitigated by asymmetry, not eliminated.} We claim
mint-after-cutoff, not contamination-proofing. A public generator can synthesize
in-distribution data, and the public \texttt{dev} split will itself be absorbed into
future training sets over time. That is exactly why the scored \texttt{test} split is
minted from a private, never-committed seed and why we report graded A/B/C novelty tiers
and embed a canary (\Cref{sec:contam}): the durable contribution is the protocol --- a
fresh draw from a private seed --- not a permanently clean file. Two honest caveats
follow: the tier gap is informative only \emph{if} a study actually observes it, and the
held-out seed and Tier B/C structures are deliberately withheld, which is a freshness
trade-off, not an oversight.

\paragraph{Dual use.} \tool{} is, literally, a bad-code generator: its transforms overlap
with source obfuscation, and a tool that emits deliberately unmaintainable code invites
the obvious misuse concern. Two properties bound that risk. First, the intended use is
\emph{research} --- building contamination-controlled, equivalence-verified, labelled data
for evaluating and probing code models --- and the generator changes presentation at zero
semantic cost; it does not introduce vulnerabilities or alter behavior. Second, and more
to the point, every output is \emph{clearly labelled} (each anti-pattern is annotated in
the emitted source) and \emph{oracle-verified}, which makes the output a poor instrument
for passing degraded code off as legitimate: the messiness is documented by construction.
We judge the research value --- a class of controlled code data that mining cannot
supply --- to outweigh a dual-use surface that obfuscation toolkits already occupy, and
note that the artifact ships under an open license with this discussion attached.
